\chapter{Future Work}
\label{ch:future-work}

This chapter delineates the roadmap for transitioning the Wolverine platform from an experimental prototype to a production-grade cyber-intelligence research framework. Grounded in the limitations identified in \Cref{ch:evaluation,ch:discussion}, the proposed advancements are organized across near-term engineering refinements, analytical and algorithmic enhancements, dynamic evaluation frameworks, architectural deployment evolution, and long-term scientific research directions.

\section{Near-Term Engineering Refinements}
\label{sec:future-engineering}
Immediate engineering efforts must focus on hardening prototype subsystems and resolving documented architectural gaps:
\begin{itemize}
    \item \textbf{Durable Replay Protection:} Replacing the in-memory \texttt{Set<string>} deduplication cache in WolverineDB with a persistent Redis or PostgreSQL-backed transactional ledger to guarantee replay protection across process restarts.
    \item \textbf{Physical BFT Network Transport:} Migrating \texttt{wolverine-db} from in-process \texttt{DirectMemoryNetworkTransport} to real TCP/gRPC socket connections with TLS mutual authentication to validate Byzantine fault tolerance across physical network boundaries.
    \item \textbf{Hardware Key Management (HSM/KMS):} Replacing simulated software key derivation with PKCS\#11 interfaces to genuine Hardware Security Modules or cloud KMS services (e.g., AWS KMS, HashiCorp Vault).
    \item \textbf{Persistent Blockchain Anchoring:} Connecting the Merkle root checkpoint daemon (\texttt{src/anchors/}) to a live EVM testnet via Web3/Ethers.js RPC clients rather than in-memory hash maps.
\end{itemize}

\section{Analytical and Algorithmic Enhancements}
\label{sec:future-analytics}
To transcend classical heuristic string matching, future research should integrate modern machine learning and semantic entity resolution:
\begin{itemize}
    \item \textbf{Learned Entity Resolution:} Replacing static linear weights with supervised gradient-boosted decision trees (e.g., XGBoost) or graph neural networks (GNNs) trained on synthetic pair annotations.
    \item \textbf{Dynamic Behavioral Overlap Calculation:} Replacing the hardcoded $s_{\text{act}} = 0.5$ baseline with dynamic Jaccard or cosine similarity over normalized posting timestamps, transaction frequency vectors, and active diurnal time zones.
    \item \textbf{Stylometric and Semantic Embeddings:} Augmenting the 4-character prefix cue ($s_{\text{cue}}$) with pre-trained transformer embeddings (e.g., sentence-transformers) to detect stylistic and syntactic similarities across forum posts while preserving privacy.
    \item \textbf{Probabilistic Graph Traversals:} Extending graph projection beyond depth $d = 4$ BFS by implementing personalized PageRank and community detection algorithms (e.g., Louvain clustering).
\end{itemize}

\section{Dynamic Evaluation Frameworks}
\label{sec:future-evaluation}
To replace mock evaluation outputs with rigorous empirical benchmarks, the following experimental setups must be executed:
\begin{itemize}
    \item \textbf{Dynamic Full-Dataset Benchmarking:} Implementing an automated evaluation pipeline that computes dynamic confusion matrices ($TP, FP, FN$), Precision-Recall curves, and $F_\beta$ scores across all 89,605 accounts in the synthetic dataset under parameterized noise levels (0\% to 25\%).
    \item \textbf{Ablation Studies:} Measuring the individual performance contribution of each feature weight ($w_a, w_d, w_t, w_o, w_c$) by iteratively removing features and evaluating ROC-AUC degradation.
    \item \textbf{Human Analyst User Trials:} Conducting double-blind laboratory experiments measuring investigative time-to-insight and error rates between analysts utilizing the Wolverine platform versus traditional manual query interfaces.
\end{itemize}

\section{Deployment and Infrastructure Evolution}
\label{sec:future-deployment}
To achieve horizontal scalability, the infrastructure must evolve beyond single-host Docker Compose:
\begin{itemize}
    \item \textbf{Distributed Graph Persistence:} Migrating from in-memory adjacency maps to dedicated distributed graph databases (e.g., Neo4j, Apache AGE, or optimized PostgreSQL recursive CTEs) with indexing on composite entity IDs.
    \item \textbf{Kubernetes Orchestration:} Defining Helm charts and Kubernetes manifests with resource quotas, horizontal pod autoscaling (HPA), and automated volume replication across worker nodes.
    \item \textbf{Asynchronous Event-Driven Ingestion:} Expanding the Redis stream outbox architecture to Apache Kafka to support high-throughput streaming ingestion from external crawl engines.
\end{itemize}

\section{Governance, Explainability, and Ethical Evolution}
\label{sec:future-governance}
Future iterations must incorporate advanced algorithmic governance:
\begin{itemize}
    \item \textbf{Explainable Entity Linkage:} Generating human-interpretable confidence decompositions for every resolved edge (e.g., providing feature-by-feature SHAP value breakdowns in the Analyst UI).
    \item \textbf{Automated Provenance Auditing:} Integrating cryptographic Merkle proof generators directly into the user interface, enabling analysts to export one-click tamper-evident forensic dossiers for judicial review.
    \item \textbf{Bias and Disparate Impact Auditing:} Establishing regular algorithmic audits to detect and mitigate potential systematic misattribution biases across synthetic demographic personas.
\end{itemize}

\section{Long-Term Scientific Research Directions}
\label{sec:future-research}
The prototype opens several foundational research questions at the intersection of cryptography, machine learning, and cyber-intelligence:
\begin{enumerate}
    \item \textit{Can zero-knowledge proofs (ZKPs) allow intelligence agencies to verify cross-border entity correlations without revealing underlying proprietary collection data?}
    \item \textit{How do adversarial evasion strategies (e.g., automated handle shuffling, sybil account flooding) degrade heuristic vs. learned entity resolution under real-world domain shifts?}
    \item \textit{What formal mathematical bounds can be established for prompt-injection resistance in retrieval-augmented LLM architectures ingesting hostile external text?}
\end{enumerate}
